\chapter*{Danh mục từ viết tắt và thuật ngữ}
\addcontentsline{toc}{chapter}{Danh mục từ viết tắt và thuật ngữ}

\begin{xltabular}{\textwidth}{lp{4.2cm}X}
\toprule
\textbf{Từ viết tắt} & \textbf{Thuật ngữ Tiếng Anh} & \textbf{Diễn giải Tiếng Việt} \\
\midrule
\endfirsthead
\toprule
\textbf{Từ viết tắt} & \textbf{Thuật ngữ Tiếng Anh} & \textbf{Diễn giải Tiếng Việt} \\
\midrule
\endhead
\bottomrule
\endfoot
\bottomrule
\endlastfoot

AAA & Authentication, Authorization, and Accounting & Hệ thống xác thực, cấp quyền và lưu vết tài khoản \\
AP & Access Point & Điểm truy cập không dây (Trạm phát Wi-Fi) \\
BSS & Basic Service Set & Tập dịch vụ cơ bản (Mô hình mạng không dây có AP) \\
BSS Coloring & Basic Service Set Coloring & Kỹ thuật tô màu định danh bản tin trong Wi-Fi 6 \\
BYOD & Bring Your Own Device & Mô hình người dùng tự mang thiết bị cá nhân \\
CAPWAP & Control and Provisioning of Wireless Access Points & Giao thức điều khiển và cấp phép cho điểm truy cập \\
CSMA/CA & Carrier Sense Multiple Access with Collision Avoidance & Đa truy nhập cảm nhận sóng mang tránh xung đột \\
CSMA/CD & Carrier Sense Multiple Access with Collision Detection & Đa truy nhập cảm nhận sóng mang phát hiện xung đột \\
DCA & Dynamic Channel Assignment & Thuật toán gán kênh tần số tự động \\
DHCP & Dynamic Host Configuration Protocol & Giao thức cấu hình và cấp phát địa chỉ IP tự động \\
DSSS & Direct Sequence Spread Spectrum & Trải phổ chuỗi trực tiếp \\
DTHU & Dong Thap University & Trường Đại học Đồng Tháp \\
EAP & Extensible Authentication Protocol & Giao thức xác thực mở rộng \\
ESS & Extended Service Set & Tập dịch vụ mở rộng (Hệ thống liên kết nhiều BSS) \\
GVHD & Giảng viên hướng dẫn & Cán bộ hướng dẫn khoa học của sinh viên \\
IBSS & Independent Basic Service Set & Tập dịch vụ cơ bản độc lập (Mô hình mạng Ad-hoc) \\
IEEE & Institute of Electrical and Electronics Engineers & Viện Kỹ sư Điện và Điện tử \\
IP & Internet Protocol & Giao thức mạng Internet / Địa chỉ mạng IP \\
MAC & Medium Access Control & Lớp điều khiển truy nhập môi trường truyền dẫn \\
MIMO & Multiple Input Multiple Output & Công nghệ đa anten phát và đa anten thu \\
MU-MIMO & Multi-User Multiple Input Multiple Output & Công nghệ MIMO đa người dùng \\
OFDMA & Orthogonal Frequency Division Multiple Access & Đa truy nhập phân chia theo tần số trực giao \\
PBR & Policy-Based Routing & Định tuyến theo chính sách \\
PHY & Physical Layer & Lớp vật lý trong mô hình mạng \\
PoE / PoE+ & Power over Ethernet & Công nghệ cấp nguồn điện qua cáp mạng mạng LAN \\
QoS & Quality of Service & Chất lượng dịch vụ mạng \\
RADIUS & Remote Authentication Dial-In User Service & Dịch vụ máy chủ xác thực người dùng tập trung \\
RRM & Radio Resource Management & Cơ chế quản lý tài nguyên sóng vô tuyến tự động \\
RSSI & Received Signal Strength Indicator & Chỉ số đo lường cường độ tín hiệu sóng thu được \\
RTS/CTS & Request To Send / Clear To Send & Cơ chế gửi yêu cầu và nhận xác nhận truyền tin \\
RU & Resource Unit & Đơn vị tài nguyên tần số trong chuẩn Wi-Fi 6 \\
SNR & Signal-to-Noise Ratio & Tỷ số tín hiệu trên tạp âm (tỷ số tín hiệu/nhiễu) \\
SSID & Service Set Identifier & Tên định danh của mạng không dây \\
STA & Station & Thiết bị máy trạm người dùng đầu cuối \\
TPC & Transmit Power Control & Cơ chế tự động điều chỉnh công suất phát sóng \\
VLAN & Virtual Local Area Network & Mạng cục bộ ảo logic \\
WIDS/WIPS & Wireless Intrusion Detection / Prevention System & Hệ thống phát hiện và ngăn chặn xâm nhập không dây \\
WLAN & Wireless Local Area Network & Mạng cục bộ không dây \\
WLC & Wireless LAN Controller & Bộ điều khiển trung tâm quản lý mạng không dây \\
WPA / WPA3 & Wi-Fi Protected Access & Các chuẩn giao thức bảo mật và mã hóa mạng Wi-Fi \\

\end{xltabular}
